% ============================================================
% Secao 9: Loss Composta
% ============================================================

\section{Loss Composta}
\label{sec:loss}

O \aletheion{} utiliza uma loss composta com 14 componentes, agrupadas em
cross-entropy padr\~ao e 13 losses epist\^emicas com \textit{annealing}
progressivo.

\subsection{Formula\c{c}\~ao Geral}

\begin{equation}
    \calL_{\text{total}} = \lambda_{\text{CE}} \cdot \calL_{\text{CE}} + \alpha(t) \cdot \sum_{i=1}^{13} \lambda_i \cdot \calL_i
    \label{eq:total_loss}
\end{equation}

onde $\alpha(t) \in [0, 1]$ \'e o fator de \textit{annealing} e $\lambda_i$
s\~ao pesos por componente.

\subsection{Cross-Entropy}

A loss principal \'e a cross-entropy padr\~ao para modelagem autorregressiva:

\begin{equation}
    \calL_{\text{CE}} = -\frac{1}{|\calS|} \sum_{t \in \calS} \log P(x_t \mid x_{<t})
\end{equation}

onde $P(x_t \mid x_{<t}) = \text{softmax}(\bm{l}_t / \tau)_{x_t}$ e
$\calS$ exclui tokens de padding.

\subsection{12 Losses Epist\^emicas}

\begin{table}[H]
\centering
\caption{Componentes da loss composta com pesos default.}
\label{tab:losses}
\begin{tabular}{clcp{6cm}}
\toprule
\textbf{\#} & \textbf{Nome} & \textbf{$\lambda$} & \textbf{Objetivo} \\
\midrule
1 & VARO & 0.10 & Calibra $q_1 \cdot q_2$ vs acur\'acia \\
2 & VI & 0.01 & Mant\'em $\phi \geq \phi_{\text{cr\'itico}}$ + penalidade severidade \\
3 & MAD Calibra\c{c}\~ao & 0.05 & Calibra confian\c{c}a vs acur\'acia \\
4 & M\'etrica & 0.001 & $\kappa(\bG) \to 1$ (condition number) \\
5 & Eidos & 0.005 & Balanceamento de eixos \\
6 & Conflito & 0.005 & Minimiza $\gamma^2$ \\
7 & Consci\^encia & 0.003 & Energia $\geq 0.2$ \\
8 & Grounding & 0.005 & Entropia tarefas + calibra ambiguidade \\
9 & Plasticidade & 0.002 & Plasticidade $\geq 0.3$ \\
10 & Fronteira & 0.002 & Maximiza explora\c{c}\~ao \\
11 & MOPsi & 0.003 & Alinhamento $\psi$-confian\c{c}a \\
12 & Contrastivo & 0.003 & Anti-colapso das vis\~oes \\
\bottomrule
\end{tabular}
\end{table}

\subsection{Detalhamento das Losses}

\paragraph{VARO (Calibra\c{c}\~ao).}
Calibra o sinal combinado de incerteza contra a acur\'acia real:

\begin{equation}
    \calL_{\text{VARO}} = \E\big[(q_1 \cdot q_2 - p_{\text{correto}})^2\big]
\end{equation}

onde $p_{\text{correto}} = P(x_t = x_t^* \mid x_{<t})$ \'e a probabilidade
softmax do token correto. Isso for\c{c}a o produto $q_1 \cdot q_2$ a estimar
a confian\c{c}a real do modelo em cada predi\c{c}\~ao.

\paragraph{Regulariza\c{c}\~ao VI.}
Penaliza manifold degradado e severidade excessiva (\Cref{sec:vi}):

\begin{equation}
    \calL_{\text{VI}} = \E\big[\text{ReLU}(\phi_{\text{cr\'itico}} - \phi)^2\big] + 0.3 \cdot \E[s^2], \quad \phi_{\text{cr\'itico}} = 0.5
\end{equation}

\paragraph{MAD Calibra\c{c}\~ao.}
BCE entre confian\c{c}a e acur\'acia (\Cref{eq:mad_loss}).

\paragraph{Regulariza\c{c}\~ao M\'etrica.}
\begin{equation}
    \calL_{\text{metric}} = \log(\kappa(\bG) + 1)
\end{equation}

\paragraph{Eidos.}
Penaliza desbalanceamento de eixos:

\begin{equation}
    \calL_{\text{eidos}} = \E\big[(\text{std}(\text{axis\_balance}) - 0.15)^2\big]
\end{equation}

\paragraph{Conflito.}
Minimiza intensidade de conflito $\phi$-$\psi$:

\begin{equation}
    \calL_{\text{conflito}} = \E[\gamma^2]
\end{equation}

\paragraph{Consci\^encia.}
Impede colapso de energia:

\begin{equation}
    \calL_{\text{cons}} = \E\big[\text{ReLU}(0.2 - \text{energy})^2\big]
\end{equation}

\paragraph{Grounding.}
Entropia das probabilidades de tarefa + calibra\c{c}\~ao de ambiguidade:

\begin{equation}
    \calL_{\text{ground}} = -\E\Big[\sum_k p_k \log p_k\Big] + \E\big[(\text{ambiguity} - q_1)^2\big]
\end{equation}

\paragraph{Plasticidade.}
Impede deple\c{c}\~ao total:

\begin{equation}
    \calL_{\text{plast}} = \E\big[\text{ReLU}(0.3 - \text{plast})^2\big]
\end{equation}

\paragraph{Fronteira.}
Maximiza explora\c{c}\~ao de regi\~oes novas:

\begin{equation}
    \calL_{\text{frontier}} = -\E\big[f_{\text{frontier}} \cdot \log(\text{novelty} + \epsilon)\big]
\end{equation}

\paragraph{MOPsi.}
Alinha $\psi$ com confian\c{c}a:

\begin{equation}
    \calL_{\text{mopsi}} = \E\big[(\psi_{\text{score}} - \kappa)^2\big]
\end{equation}

\paragraph{Contrastivo.}
Anti-colapso: garante que as duas vis\~oes n\~ao sejam id\^enticas:

\begin{equation}
    \calL_{\text{contr}} = -\E[\log(\text{div} + \epsilon)] + 0.1 \cdot \E\big[\text{ReLU}(\text{div} - 0.8)^2\big]
\end{equation}

O primeiro termo maximiza diverg\^encia (anti-colapso), o segundo
impede diverg\^encia excessiva.

\subsection{Schedule de Annealing}
\label{sec:annealing}

As losses epist\^emicas s\~ao introduzidas gradualmente para estabilizar
o treinamento:

\begin{equation}
    \alpha(t) = \begin{cases}
        0 & \text{se } t < 0.1 \cdot T_{\text{total}} \\
        \frac{t - 0.1T}{0.4T} & \text{se } 0.1T \leq t < 0.5T \\
        1 & \text{se } t \geq 0.5T
    \end{cases}
    \label{eq:annealing}
\end{equation}

\begin{itemize}
    \item \textbf{Fase 1} (0--10\%): Apenas $\calL_{\text{CE}}$. O modelo
    aprende a distribuir tokens antes de qualquer sinal epist\^emico.

    \item \textbf{Fase 2} (10--50\%): $\alpha$ cresce linearmente de 0 a 1.
    As losses epist\^emicas s\~ao introduzidas suavemente.

    \item \textbf{Fase 3} (50--100\%): $\alpha = 1$. Todas as losses ativas
    com pesos plenos.
\end{itemize}

\begin{theorem}[Estabilidade do Annealing]
O schedule de annealing garante que os gradientes epist\^emicos n\~ao
dominam os gradientes de CE nas fases iniciais:

\begin{equation}
    \frac{\|\nabla_\theta \alpha \calL_{\text{epist}}\|}{\|\nabla_\theta \calL_{\text{CE}}\|} \leq \alpha(t) \cdot \frac{\sum_i \lambda_i \|\nabla_\theta \calL_i\|}{\|\nabla_\theta \calL_{\text{CE}}\|}
\end{equation}

Para $\alpha(t) \ll 1$ nas fases iniciais, o r\'acio \'e controlado.
\end{theorem}

\subsection{Pesos dos Componentes}

Os pesos $\lambda_i$ foram calibrados empiricamente com os seguintes
crit\'erios:

\begin{enumerate}
    \item $\lambda_{\text{CE}} = 1.0$ (ancora)
    \item Losses de regulariza\c{c}\~ao ($\leq 0.01$): impedem degenera\c{c}\~ao
    sem dominar
    \item Losses de calibra\c{c}\~ao ($\leq 0.005$): ajuste fino de sinais
    \item Losses de explora\c{c}\~ao ($\leq 0.003$): incentivam diversidade
\end{enumerate}
